% ============================================================
% Oppgave 5
% ============================================================
\hypertarget{oppgave5}{}
\section*{Oppgave 5}
\addcontentsline{toc}{section}{Oppgave 5}
Datasekvensen er:
\[
x[n]
=
\{10,14,2,6,8,20,1,17,14,7,15,16\}
\]
\subsection*{a)}
Filteret skal beregne middelverdien av tre etterfølgende
ingangsverdier.
Den ikke-rekursive løsningen blir derfor:
\[
\boxed{
y[n]
=
\frac{x[n]+x[n-1]+x[n-2]}{3}
}
\]
To tidligere inngangsverdier må lagres. Verdier før starten av
sekvensen settes lik null.
Den rekursive løsningen kan finnes ved å skrive:
\[
y[n-1]
=
\frac{x[n-1]+x[n-2]+x[n-3]}{3}
\]
Differansen mellom $y[n]$ og $y[n-1]$ blir:
\[
y[n]-y[n-1]
=
\frac{x[n]-x[n-3]}{3}
\]
Dermed får vi:
\[
\boxed{
y[n]
=
y[n-1]
+
\frac{x[n]-x[n-3]}{3}
}
\]
Den ikke-rekursive og den rekursive løsningen gir samme
matematiske resultat når beregningene utføres med uendelig
presisjon.
For den oppgitte sekvensen, med tre nuller før starten og to nuller
etter slutten, blir utsignalene:
\[
\begin{aligned}
y[n]=\{&
0,\,
0,\,
0,\,
3.3333,\,
8.0000,\,
8.6667,\,
7.3333,\,
5.3333,\,
11.3333,\,
9.6667,\,
12.6667,\\
&
10.6667,\,
12.6667,\,
12.0000,\,
12.6667,\,
10.3333,\,
5.3333
\}
\end{aligned}
\]
og den tilsvarende rekursive løsningen er:
\[
\begin{aligned}
yr[n]=\{&
0,\,
0,\,
0,\,
3.3333,\,
8.0000,\,
8.6667,\,
7.3333,\,
5.3333,\,
11.3333,\,
9.6667,\,
12.6667,\\
&
10.6667,\,
12.6667,\,
12.0000,\,
12.6667,\,
10.3333,\,
5.3333
\}
\end{aligned}
\]
Differansen mellom de to utgangene er:
\[
\begin{aligned}
d[n]=y[n]-yr[n]=\{&
0,\,
0,\,
0,\,
0,\,
0,\,
0,\,
0,\,
0,\,
1.7764\times 10^{-15},\,
0,\,
0,\,
0,\,
0,\\
&
0,\,
0,\,
0,\,
1.7764\times 10^{-15},\,
8.8818\times 10^{-16}
\}
\end{aligned}
\]
Maksimal absolutt differanse er:
\[
\max |d[n]| = 1.7764\times 10^{-15}
\]
Octave-koden for de to løsningene er:
\begin{verbatim}
xx = [10,14,2,6,8,20,1,17,14,7,15,16];
x = [zeros(1,3), xx, zeros(1,2)];
y = zeros(size(x));
for n = 3:length(x)
    y(n) = (x(n)+x(n-1)+x(n-2))/3;
end
yr = zeros(size(x));
for n = 4:length(x)
    yr(n) = yr(n-1) + (x(n)-x(n-3))/3;
end
d = y - yr;
maxAbsDiff = max(abs(d));
\end{verbatim}
\subsection*{b)}
Tallet:
\[
\frac{1}{3}
=
0.3333\ldots
\]
kan ikke representeres eksakt med et endelig antall binære
desimaler.
De oppgitte tilnærmingene er:
\[
\begin{array}{c|c|c}
\text{Binære desimaler} &
\text{Tilnærming} &
\text{Feil fra }1/3 \\
\hline
3 & 0.375000 & +0.041667 \\
4 & 0.312500 & -0.020833 \\
5 & 0.343750 & +0.010417 \\
6 & 0.328125 & -0.005208
\end{array}
\]
Feilen blir mindre når antall tilgjengelige binære desimaler
øker.
For hver verdi av $c$ testes tre realiseringer:
\[
\begin{aligned}
y_{\mathrm{sum}}[n]&=c(x[n]+x[n-1]+x[n-2]),\\
y_{\mathrm{prod}}[n]&=cx[n]+cx[n-1]+cx[n-2],\\
y_r[n]&=y_r[n-1]+c(x[n]-x[n-3]).
\end{aligned}
\]
Følgende Octave-kode utfører testen. Referansen \texttt{ref} bruker
den eksakte verdien $1/3$ med Octaves flyttallspresisjon:
\begin{verbatim}
cs = [0.375, 0.3125, 0.34375, 0.328125];
ref = zeros(size(x));
for n = 3:length(x)
    ref(n) = (x(n)+x(n-1)+x(n-2))/3;
end
for q = 1:length(cs)
    c = cs(q);
    ysum = zeros(size(x));
    yprod = zeros(size(x));
    yr = zeros(size(x));
    for n = 3:length(x)
        ysum(n) = c*(x(n)+x(n-1)+x(n-2));
        yprod(n) = c*x(n)+c*x(n-1)+c*x(n-2);
    end
    for n = 4:length(x)
        yr(n) = yr(n-1)+c*(x(n)-x(n-3));
    end
    printf("c = %.6f: %.12g %.12g %.12g %.12g\n", c, ...
        max(abs(ysum-ref)), max(abs(yr-ref)), ...
        max(abs(ysum-yr)), max(abs(ysum-yprod)));
end
\end{verbatim}
Resultatene, der hver verdi er største absolutte avvik i hele
utgangssekvensen, blir:
\[
\begin{array}{c|c|c|c}
\text{Binære desimaler} & c &
\max|y_{\mathrm{sum}}-\mathrm{ref}| &
\max|y_r-\mathrm{ref}| \\
\hline
3 & 0.375000 & 1.583333 & 1.583333 \\
4 & 0.312500 & 0.791667 & 0.791667 \\
5 & 0.343750 & 0.395833 & 0.395833 \\
6 & 0.328125 & 0.197917 & 0.197917
\end{array}
\]
For alle fire verdiene av $c$ får vi dessuten:
\[
\max|y_{\mathrm{sum}}-y_r|=0
\qquad \text{og} \qquad
\max|y_{\mathrm{sum}}-y_{\mathrm{prod}}|=0.
\]
Feilen mot referansen halveres for hver ekstra binære desimal. I
denne testen er de tre realiseringene like, fordi de oppgitte
verdiene av $c$ kan representeres eksakt binært, inngangsverdiene
er heltall, og Octave bruker høy presisjon i mellomregningene.
En virkelig prosessor med begrenset antall bits kan i tillegg runde
av hver mellomregning. Da kan den rekursive strukturen videreføre
avrundingsfeil, mens $y_{\mathrm{prod}}$ kan få flere
avrundingsfeil enn $y_{\mathrm{sum}}$ fordi den utfører tre
multiplikasjoner i stedet for en.